\documentclass[12pt]{article}
\usepackage[utf8]{inputenc}
\usepackage[T1]{fontenc}
\usepackage{hyperref}
\title{Makhuwa Language Resources for K-12 Curriculum Development}
\author{Compiled by Claude Sonnet 5 (claude-sonnet-5) for Andr\'as Kornai}
\date{}
\begin{document}
\maketitle

\section{Introduction}
Makhuwa (ISO 639-3 vmw; Glottolog makh1264) is a Bantu language of northern Mozambique, with millions of speakers. The searches found few directly usable school resources.

\section{Primary Educational Resources}
\begin{itemize}
\item University of S\~ao Paulo doctoral thesis on Emakhuwa in bilingual education, with proposed classroom exercises: \url{https://teses.usp.br/teses/disponiveis/8/8142/tde-10112017-120739/es.php}
\end{itemize}

\section{Linguistic and Reference Materials}
\begin{itemize}
\item Wikipedia overview: \url{https://en.wikipedia.org/wiki/Makhuwa_language}
\item Historical handbook, Maples, \emph{Collections for a handbook of the Makua language}, held by SOAS: \url{https://digital.soas.ac.uk/LOAA000003/00001/citation}
\item Cambridge talk abstract on Makhuwa structure, by Jenneke van der Wal: \url{https://talks.cam.ac.uk/talk/index/40041/}
\end{itemize}

\section{Religious and Community Resources}
None were found in the searches.

\section{Digital and Audio Resources}
None were found. Software localization pages for Makhuwa varieties appeared, but they were not checked and are not cited.

\section{Curriculum Development Context}
The thesis above reports that Portuguese is the usual language of instruction. No modern dictionary, grammar or textbook series was found, so those would need targeted follow-up searches. They should include Mozambican bilingual education programmes and the Makhuwa-language Bible and literacy materials.
\end{document}